\documentclass[11pt,a4paper]{article}
\usepackage[T1]{fontenc}
\usepackage{lmodern,amsmath,amssymb,amsthm,mathtools}
\usepackage[margin=25mm,headheight=15pt]{geometry}
\usepackage{booktabs,array,tabularx}
\usepackage[dvipsnames]{xcolor}
\usepackage{microtype,enumitem,fancyhdr,needspace}
\usepackage[colorlinks=true,linkcolor=MidnightBlue,citecolor=MidnightBlue,urlcolor=MidnightBlue]{hyperref}
\usepackage{bookmark}
\hypersetup{pdftitle={Native Stepping and Exchange Cadence in the Reciprocal Chiral Exchange Candidate},
pdfauthor={Hilmir Frímann Halldórsson (pzychozen)},
pdfsubject={Finite-step accounting, initial receiving, and local capacity. Research manuscript v0.1. Accepted mathematics; manuscript review pending. Prepared 9 November 2026.}}
\setlength{\parskip}{4pt}\setlength{\parindent}{0pt}
\setlength{\emergencystretch}{3em}
\pagestyle{fancy}\fancyhf{}
\fancyhead[L]{\small Native stepping and exchange cadence}
\fancyhead[R]{\small v0.1}
\fancyfoot[C]{\thepage}\renewcommand{\headrulewidth}{0.3pt}
\numberwithin{equation}{section}
\newtheorem{theorem}{Theorem}[section]
\newtheorem{proposition}[theorem]{Proposition}
\newcommand{\R}{\mathbb R}\newcommand{\C}{\mathbb C}
\newcommand{\B}{B}\newcommand{\dd}{\,\mathrm d}
\newcommand{\norm}[1]{\left\|#1\right\|}
\newcolumntype{Y}{>{\raggedright\arraybackslash}X}
\begin{document}
\hypersetup{pageanchor=false}
\begin{titlepage}
\vspace*{2mm}
{\small\color{MidnightBlue} TRI-OCTAGON MATHEMATICAL RESEARCH\par}
\vspace{7mm}
{\LARGE\bfseries\raggedright Native Stepping and Exchange Cadence\\
in the Reciprocal Chiral Exchange\\Candidate\par}
\vspace{5mm}
{\large Finite-step accounting, initial receiving,\\and local capacity\par}
\vspace{6mm}
{\large Hilmir Fr\'imann Halld\'orsson (pzychozen)\par}
\vspace{3mm}
{Research manuscript v0.1\quad 9 November 2026\par}
{\itshape Accepted mathematics; manuscript review pending.\par}
\vspace{4mm}
\begin{abstract}
The conservative reciprocal chiral exchange candidate admits an exact finite native update between exchange intervals. We examine the inherited native-first composition with a fixed symbolic positive cadence. Native baseline increments must be subtracted when attributing recovery to exchange; the invariant storage floor can change at a native jump, so the old capacity cannot in general be carried across unchanged. For the coherent empty launch $(1,e^{i\phi},1)$, $0<\phi\le1/100$ radians, the exact first baseline has both signs and the native update leaves the precise normalized launch family. Nevertheless, its synchronized source remains nonzero throughout this interval. For each fixed positive phase and finite positive exchange tuple, the actual post-update trajectory has strict initial native debit and receiver gain on a sufficiently short interval before the next jump. This interval is pointwise in phase and parameters. At fixed nonzero native detuning, a local capacity calculation gives a new-to-old algebraic ceiling ratio whose limit is approximately $1.132216337515111$; no attained-transfer claim follows. Strong continuity gives a one-step small-cadence limit equal to the finite native update, with a separate qualified fixed-number-of-steps limit. The later capacity-sized receiving-and-recovery episode under native stepping remains unproved, and no physical cadence is selected.
\end{abstract}
\vfill
{\small Research author and geometric direction: Hilmir Fr\'imann Halld\'orsson.
GPT formulated and internally reviewed the accepted investigation.
Codex independently investigated the result and prepared this manuscript.
The new manuscript and proposed publication payload await GPT/author review.
No external peer review or machine-formal verification is claimed.}
\end{titlepage}
\setcounter{page}{2}\hypersetup{pageanchor=true}
\section{Introduction}\label{sec:introduction}
The reciprocal chiral exchange candidate couples a finite native state, a periodic carrier and a conservative receiver. Its continuous exchange flow has an exact storage account, general projection ceilings and a receiving-and-recovery certificate for a particular coherent empty preparation~\cite{rx,pc,rr}. The native model also has an exact discrete map. Inserting that map raises a separate question: which changes in native storage remain attributable to exchange, and what can be inferred about the interval between native updates?

We keep the already defined native-first hybrid~\cite[Section 6]{rx}. The native map acts once, and the exact exchange flow then acts for a fixed symbolic interval $h>0$. The native coefficients do not scale with $h$. Carrier and receiver states are retained at each jump. Neither a physical clock nor an alternative ordering is supplied.

The principal positive result is Theorem~\ref{thm:initial}: the actual synchronized post-native state has nonzero drive for every phase in the stated launch interval and therefore initially supplies the carrier and receiver. Its short-time interval may depend on phase and the fixed finite coefficients. The exact accounts in Section~\ref{sec:accounts}, the first-jump calculation in Section~\ref{sec:firstjump}, the local ceiling comparison in Section~\ref{sec:capacity} and the limits in Section~\ref{sec:cadence} explain what this initial result does and does not establish.

The global exchange theorem, the native map and hybrid ordering, the exchange invariants, and the restricted retained-receiver ratio are inherited. Their indexed use, the explicit first-jump specialization, the full-phase synchronized source estimate, the local capacity calculation and the qualified limit deductions are the accepted present investigation~\cite{investigation,acceptance}. In particular, failure to satisfy the older launch hypothesis is not a proof that later hybrid receiving or recovery is impossible. It means that a new estimate would be required.

\subsection{Notation and conventions}\label{sec:notation}
Native amplitudes retain the published normalized convention, and angles are in radians. The native coefficients $\epsilon,g,\lambda,k_i$ are distinct from the carrier width $\epsilon_c$, yaw $\delta_g$, readout gain $G$ and operator norm $\|K\|$. The full post-jump vector is $f_{\rm post}$, whereas $f_{{\rm raw},\pm}$ denotes an individual raw port. The cadence is $h$; $h_e$ is a kernel-projection quantity. The carrier length is $L_c$ and the invariant floor is $\mathcal L$.

Relative to the accepted proof, we write $\mathcal L$ for its floor $L$, $\mathbf r_0$ for its zero-phase vector $r$, $\xi$ for its phase-derivative vector $t$, $\zeta$ for its unknown second-order real vector $v_2$, and $q_{\rm cap}$ for its local capacity coefficient $Q$. These changes are one-to-one. Carrier $r_c$, exchange time $t$, prestage $v$ and accumulated work $Q_{A\to B}$ retain their separate meanings. The complete equation and source correspondence accompanies the manuscript.

\section{Exact map and exchange framework}\label{sec:model}
\subsection{Simultaneous native update}\label{sec:native}
For $\Omega=a+ib\in\C^3$ write $x=(a,b)\in\R^6$ and set $I=|\Omega|^2$. The inherited parameter tuple is~\cite[Section 6, native tuple]{rx}
\begin{equation}\label{eq:a01}% accepted equation 1
\begin{gathered}
\epsilon=\frac1{10},\qquad g=\frac1{60},\qquad
\lambda=\frac1{300},\qquad
(k_1,k_2,k_3)=\left(\frac{999}{1000},1,\frac{1001}{1000}\right).
\end{gathered}
\end{equation}
Every component of the prestage uses the same original state:
\begin{equation}\label{eq:a02}% accepted equation 2
\begin{gathered}
d_{{\rm pre},i}=\epsilon(k_i-|\Omega_i|^2)\Omega_i+
g\sum_{j\ne i}(\Omega_j-\Omega_i),\qquad v=\Omega+d_{\rm pre}.
\end{gathered}
\end{equation}
Every synchronization increment then uses that same prestage:
\begin{equation}\label{eq:a03}% accepted equation 3
\begin{gathered}
\delta_i=\lambda\sum_{j\ne i}\sin(3[\operatorname{Arg}_0v_j-\operatorname{Arg}_0v_i]),
\qquad F_i(\Omega)=e^{i\delta_i}v_i.
\end{gathered}
\end{equation}
The convention is $\operatorname{Arg}_0(0)=0$ and exact zero output components are retained. At $\lambda=0$ the implementation bypasses phase extraction. At the present nonzero strength it applies the simultaneous harmonic-3 synchronizer. In the three-node source, both cyclic neighbours are precisely the other two nodes, and the coupling matrix is $L_3=\mathbf1\mathbf1^T-3\mathrm{Id}$~\cite{nativecode}.

The source calls $\lambda$ \texttt{phase\_strength}; its configuration does not supply numerical defaults for the tuple. The ideal expressions $e^{i\delta_i}v_i$ and $|v_i|e^{i(\operatorname{Arg}_0v_i+\delta_i)}$ agree, including at zero. Finite ideal-map statements are distinct from complex128 rounding and overflow. The prestage increment is generally not $F\Omega-\Omega$. No normalization, continuous-gradient replacement or coefficient scaling is introduced.

\subsection{Raw ports, flow and original stores}\label{sec:exchange}
Write $\alpha=\alpha_N>0$, $\beta=\beta_R>0$, $\eta=\eta_N>0$ and $\kappa=\eta_R>0$ for fixed finite exchange coefficients. At the fixed yaw $\delta_g=\pi/4$ the inherited registration is
\begin{equation}\label{eq:a04}% accepted equation 4
\begin{gathered}
\chi=a\times b,\quad f(x)=K\chi,\quad J(x)=D_xf(x),\\ K=\frac G{\sqrt7}\begin{pmatrix}\sqrt2&\sqrt2&\sqrt3\\
\sqrt2&\sqrt2&-\sqrt3\end{pmatrix},\quad G=\kappa_{\rm in}\lambda_g>0 .
\end{gathered}
\end{equation}
The raw signs are unaffected by ambient contact-tangent orientations. Retain the circle $s\in\R/L_c\mathbb Z$, finite $L_c>0$, speed $c>0$, contact centres $s_\pm=L_c/2\mp\eta_{\rm geom}$ and whole-circle shortest separation $2\eta_{\rm geom}\le L_c/2$. For $0<\epsilon_c<\eta_{\rm geom}$, the normalized profiles are
\[
b_\sigma(s)=\epsilon_c^{-1}
\left(1-\frac{|\operatorname{wrap}_{L_c}(s-s_\sigma)|}{\epsilon_c}\right)_+ .
\]
Put $\B q=b_+q_++b_-q_-$ and $\B^*p=(\int b_+p\,ds,\int b_-p\,ds)^T$, with integrals over one whole circle, and define $y=\B^*p$.
The disjoint finite-width periodic tents belong to $H^1$ and satisfy
$\B^*\B=m\mathrm{Id}_2$, $m=2/(3\epsilon_c)>0$ and $\|\B\|=\|\B^*\|=\sqrt m$.
Their first weak derivatives lie in $L^2$; their continuity prevents corner delta terms in that derivative. No second $L^2$ derivative is presumed~\cite[Section 3]{rx}.

The unchanged feedback equations are
\begin{equation}\label{eq:a05}% accepted equation 5
\begin{gathered}
\dot x=-\frac{\eta}{2\alpha}J(x)^Ty,\quad
p_t+c(r_c)_s=B(\eta f(x)-\kappa z),\quad
(r_c)_t+cp_s=0,\quad \dot z=\frac{\kappa}{\beta}y.
\end{gathered}
\end{equation}
and the original trial stores are
\begin{equation}\label{eq:a06}% accepted equation 6
\begin{gathered}
H_N=\frac\alpha2 I,\quad H_C=\frac12(\|p\|_2^2+\|r_c\|_2^2),
\quad H_R=\frac\beta2|z|^2,\quad E=H_C+H_R,\quad H_{\rm tot}=H_N+E.
\end{gathered}
\end{equation}
The one-half factors and the factor $1/(2\alpha)$ in the native equation are retained. Trial-storage commensurability supplies no new physical calibration.

The inherited global theorem~\cite[Theorem 3.1]{rx} gives a unique mild flow $\Phi_t$ on
\begin{equation}\label{eq:a07}% accepted equation 7
\begin{gathered}
\mathcal H=\mathbb R^6\times L^2_{\rm per}\times L^2_{\rm per}\times\mathbb R^2,\\ \|X\|_E^2=\alpha|x|^2+\|p\|_2^2+\|r_c\|_2^2+\beta|z|^2.
\end{gathered}
\end{equation}
Here $\mathcal H$ denotes the displayed Hilbert product; the weighted norm is equivalent to its usual product norm at fixed positive weights. The flow is strongly continuous, conserves $H_{\rm tot}$ and depends locally Lipschitz-continuously on preparation on bounded energy sets over finite times. Native and receiver coordinates are $C^1$. For periodic $H^1$ carrier data, $(p,r_c)\in C_tH^1_{\rm per}{}^2\cap C_t^1L^2{}^2$ with continuous dependence in that stronger preparation topology.

\subsection{Native-first chronology}\label{sec:hybrid}
Define $\widehat F(x,p,r_c,z)=(F_{\R}x,p,r_c,z)$ and retain
\begin{equation}\label{eq:a08}% accepted equation 8
\begin{gathered}
X_{n+1}=T_hX_n=\Phi_h\widehat F X_n,\qquad X_n^+=\widehat F X_n.
\end{gathered}
\end{equation}
$X_n$ is sampled immediately before its jump, $X_n^+$ immediately after it, and $X_{n+1}$ after exchange but before the following jump. The external coordinates do not jump. This is an exact composition of the existing map and flow, not a numerical approximation to a separately specified simultaneous native ODE.

Every finite ideal native state has a finite polynomial prestage and a defined synchronized image. The global exchange theorem then supplies the next interval, proving finite hybrid existence by induction. This includes states outside the older coherent domain but gives no uniform infinite-index bound~\cite[Section 6, limits]{rx}.

Homogeneity gives $J(x)x=2f(x)$. With periodic integration, it yields
\begin{equation}\label{eq:a09}% accepted equation 9
\begin{gathered}
P_{NC}=\eta f^Ty,\quad P_{CR}=\kappa z^Ty,\\ \dot H_N=-P_{NC},\quad\dot H_C=P_{NC}-P_{CR},\quad\dot H_R=P_{CR}.
\end{gathered}
\end{equation}
For $H^1$ preparation this follows from the weak product rule and the legitimate first derivatives. The inherited approximation and continuous-dependence argument extends the integrated accounts to mild energy-space solutions. The calculation does not infer existence merely from formal cancellation.

\section{Finite accounts and the moving floor}\label{sec:accounts}
\subsection{Baseline increments and window endpoints}\label{sec:windows}
Define the native baseline at jump $n$ and the signed exchange debit on its following interval by
\begin{equation}\label{eq:a10}% accepted equation 10
\begin{gathered}
b_n=H_N(F\Omega_n)-H_N(\Omega_n),\\ D_n=\int_0^hP_{NC}(\Phi_tX_n^+)\,dt .
\end{gathered}
\end{equation}
The jump changes only native storage; integrating the inherited accounts afterward gives
\begin{equation}\label{eq:a11}% accepted equation 11
\begin{gathered}
E_{n+1}-E_n=D_n,\quad H_{N,n+1}-H_{N,n}=b_n-D_n,\\ H_{{\rm tot},n+1}-H_{{\rm tot},n}=b_n.
\end{gathered}
\end{equation}
For $N$ completed rounds,
\begin{equation}\label{eq:a12}% accepted equation 12
\begin{gathered}
H_{\rm tot}(X_N)=H_{\rm tot}(X_0)+\sum_{n=0}^{N-1}b_n,\\ E_N-E_0=\sum_{n=0}^{N-1}D_n,\\ H_{N,N}-H_{N,0}=\sum_{n=0}^{N-1}b_n-(E_N-E_0).
\end{gathered}
\end{equation}
In particular the first native jump is included. With jump times $\tau_n=nh$, the samples are $X_n=X(\tau_n^-)$, including the prepared $X_0=X(0^-)$. The comparison from $0^-$ to $(Nh)^-$ counts $b_0,\ldots,b_{N-1}$. In contrast the comparison from $0^+$ to $(Nh)^+$ counts $b_1,\ldots,b_N$.

For an arbitrary finite window use the inherited post-jump endpoint convention~\cite[Section 7]{rx}: at a jump endpoint take the post-jump value and set
\begin{equation}\label{eq:a13}% accepted equation 13
\begin{gathered}
B_N(t_a,t_b)=\sum_{t_a<\tau_j\le t_b}b_j
\end{gathered}
\end{equation}
External stocks are continuous at the jumps. There is no impulsive carrier or receiver work. With per-port powers $P_{N,\sigma}=\eta f_\sigma y_\sigma$ and $P_{R,\sigma}=\kappa z_\sigma y_\sigma$, take positive and negative parts before summing:
\begin{equation}\label{eq:a14}% accepted equation 14
\begin{gathered}
Q_{N\to C}=\int\sum_\sigma[P_{N,\sigma}]_+,\\ Q_{C\to N}=\int\sum_\sigma[-P_{N,\sigma}]_+,\\ Q_{C\to R}=\int\sum_\sigma[P_{R,\sigma}]_+,\\ Q_{R\to C}=\int\sum_\sigma[-P_{R,\sigma}]_+.
\end{gathered}
\end{equation}
All integrals here are with respect to exchange time on the specified window. Piecewise continuity and finitely many jumps make them well defined. Opposite simultaneous port powers contribute to both directions even when their signed sum is zero. Integration between jumps and addition of baseline increments give
\begin{equation}\label{eq:a15}% accepted equation 15
\begin{aligned}
\Delta H_N&=B_N-Q_{N\to C}+Q_{C\to N},\\
\Delta H_C&=Q_{N\to C}-Q_{C\to N}-Q_{C\to R}+Q_{R\to C},\\
\Delta H_R&=Q_{C\to R}-Q_{R\to C},\\
\Delta H_N-B_N&=E(t_a)-E(t_b).
\end{aligned}
\end{equation}
The accounts do not introduce a hidden reservoir. Native baseline change is distinct from carrier or receiver work and does not make the full hybrid automatically a closed energy-conserving physical system.

\subsection{Recovery and distinct denominators}\label{sec:recovery}
Only for $E_a=E(t_a)>0$ define the exchange-attributable and raw native-recovery ratios:
\begin{equation}\label{eq:a16}% accepted equation 16
\begin{gathered}
\mathcal R_{\rm ex}=\frac{\Delta H_N-B_N}{E_a}
=1-\frac{E_b}{E_a}\le1,\\ \mathcal R_{\rm raw}=\frac{\Delta H_N}{E_a}
=\mathcal R_{\rm ex}+\frac{B_N}{E_a}.
\end{gathered}
\end{equation}
The bound uses $E_b\ge0$. Corrected recovery is nonnegative exactly when $E_b\le E_a$; no positive lower bound follows from accounting alone. A positive baseline can inflate the raw ratio without additional exchange credit, whereas a negative baseline can hide credit. At $E_a=0$ the identities remain valid but these ratios are undefined.

The inherited retained-receiver result has different hypotheses and denominator~\cite[Proposition 7.1]{rx}. If $H_C(t_a)=H_R(t_a)=0$ and $Q_{N\to C}>0$, then
\begin{equation}\label{eq:a17}% accepted equation 17
\begin{gathered}
0\le\Delta H_R=H_R(t_b)=Q_{N\to C}-Q_{C\to N}-H_C(t_b)
\le Q_{N\to C}.
\end{gathered}
\end{equation}
Thus $0\le\Delta H_R/Q_{N\to C}\le1$ survives finite native jumps with complete counting. If also $Q_{R\to C}=0$, then $\Delta H_R=Q_{C\to R}$ and the same bound applies to gross receiving. Without that condition gross work may recount back-and-forth transfer. With preload, neither nonnegative retained change nor the upper bound follows by this proof. Later intervals cannot silently be relabelled empty.

Net native recovery, retained receiver stock, gross receiving work, capacity fractions and fractions of initial native stock have distinct meanings. No tagged-cohort or physical efficiency statement follows. The published continuous-exchange return bound greater than $24/25$~\cite[Theorem 3.1]{rr} requires its actual endpoint estimates and is not imported into this hybrid.

\subsection{Projection floor with arbitrary preload}\label{sec:floor}
At the fixed yaw let $e=(1,-1,0)/\sqrt2$, so $Ke=0$. Define
\begin{equation}\label{eq:a18}% accepted equation 18
\begin{gathered}
Z_2=\Omega^T\Omega,\quad h_e=|e\cdot\Omega|^2,\quad j=I-2h_e,\\ \mathcal L(\Omega)=\frac\alpha2\max(|Z_2|,j),\\ C(\Omega)=H_N-\mathcal L=\min\{C_Z,\alpha h_e\},\\ C_Z=\frac\alpha2(I-|Z_2|).
\end{gathered}
\end{equation}
$Z_2$ uses a bilinear transpose, whereas $I$ is Hermitian intensity. The identity $I^2-|Z_2|^2=4|a\times b|^2$ gives $I\ge|Z_2|$, and $h_e\ge0$ gives $I\ge j$. Hence $0\le\mathcal L\le H_N$ and $C\ge0$; negative $j$ is allowed.

For completeness, the inherited invariants can be checked with $w=K^Ty$ and $\mu=\eta/(2\alpha)$:
\begin{equation}\label{eq:a19}% accepted equation 19
\begin{gathered}
\dot a=\mu w\times b,\quad\dot b=-\mu w\times a,\quad
\dot\Omega=-i\mu w\times\Omega.
\end{gathered}
\end{equation}
The bilinear identity $\Omega^T(w\times\Omega)=0$ makes $Z_2$ constant.
Triple products give $\dot I=-4\mu w\cdot(a\times b)$. In a right-handed basis with $e$ as third axis, $w_3=0$ and differentiation gives $\dot h_e=-2\mu w\cdot(a\times b)$; therefore $\dot j=0$~\cite[Theorem 3.1]{pc}. Native coordinates are $C^1$ even for mild carrier data, so these derivatives are valid. It is the floor $\mathcal L$, not instantaneous capacity, that stays fixed within an exchange interval.

Put $\mathcal L_n^+=\mathcal L(F\Omega_n)$ and $C_n^+=C(F\Omega_n)$. Conservation and the inherited general-preparation ceiling yield
\begin{equation}\label{eq:a20}% accepted equation 20
\begin{gathered}
0\le E(t)\le E_n+C_n^+,\qquad
-E_n\le E(t)-E_n\le C_n^+ .
\end{gathered}
\end{equation}
This keeps the actual preload $E_n$. Across the native jump, subtraction of the two defining capacities gives
\begin{equation}\label{eq:a21}% accepted equation 21
\begin{gathered}
C(F\Omega)-C(\Omega)=b(\Omega)-[\mathcal L(F\Omega)-\mathcal L(\Omega)].
\end{gathered}
\end{equation}
During exchange $C+E=H_{\rm tot}-\mathcal L$ is constant; at a jump its change is $b-\Delta\mathcal L$. Added native storage need not become additional capacity.

For $t$ after jump $n$ and before the next one, the indexed cumulative bound is
\begin{equation}\label{eq:a22}% accepted equation 22
\begin{gathered}
E(t)\le H_{\rm tot}(X_0)+\sum_{j=0}^{n}b_j-\mathcal L(F\Omega_n).
\end{gathered}
\end{equation}
The current jump is included. Endpoints must use the matching convention above. A coarser finite-index consequence is $E_N\le E_0+\sum_{n=0}^{N-1}C_n^+$. Neither estimate bounds all preloaded stock by $C_n^+$ alone or supplies an infinite-step bound. Algebraic sharpness on a common invariant level is distinct from attaining its minimum on a particular exchange trajectory~\cite[Section 3]{pc}.

\section{The exact first native jump}\label{sec:firstjump}
Specialize the first-jump comparison to
\begin{equation}\label{eq:a23}% accepted equation 23
\begin{gathered}
\Omega_\phi=(1,e^{i\phi},1),\qquad 0<\phi\le1/100,\qquad p=r_c=z=0.
\end{gathered}
\end{equation}
Put $q=e^{i\phi}$, $D=1-\cos\phi$ and $d_0=1/10000$. Since $\epsilon(k_i-1)=(-d_0,0,d_0)$, the exact prestage is
\begin{equation}\label{eq:a24}% accepted equation 24
\begin{gathered}
v_1=1-d_0+g(q-1),\\ v_2=q+2g(1-q),\\ v_3=1+d_0+g(q-1).
\end{gathered}
\end{equation}
Synchronization preserves each component modulus. For any native input, expansion of the full prestage square gives the inherited identity
\begin{equation}\label{eq:a25}% accepted equation 25
\begin{gathered}
I(F\Omega)-I(\Omega)=
2\epsilon\sum_i|\Omega_i|^2(k_i-|\Omega_i|^2)
\\ {}-2g\sum_{i<j}|\Omega_i-\Omega_j|^2+\|d_{\rm pre}\|^2.
\end{gathered}
\end{equation}
Each undirected edge is paired at its two endpoints. With $h_i=(k_i-|\Omega_i|^2)\Omega_i$, its finite square and graph identity are
\begin{equation}\label{eq:a26}% accepted equation 26
\begin{gathered}
\epsilon^2\|h\|^2+2\epsilon g\operatorname{Re}(h^\dagger L_3\Omega)
+g^2\|L_3\Omega\|^2,\\ \|L_3\Omega\|^2=3\sum_{i<j}|\Omega_i-\Omega_j|^2.
\end{gathered}
\end{equation}
The mixed term and full increment have no predetermined sign~\cite[Section 6, native intensity identity]{rx}.

\begin{proposition}[First baseline and sign change]\label{prop:firstjump}
For the preparation~\eqref{eq:a23}, the full native increment and baseline are
\begin{equation}\label{eq:a27}% accepted equation 27
\begin{gathered}
\Delta I=2d_0^2-4g(2-3g)D
=\frac1{50000000}-\frac{13}{100}D,\\ b_0=\alpha\left[\frac1{100000000}-\frac{13}{200}D\right].
\end{gathered}
\end{equation}
Both signs occur within the admitted positive-phase interval.
\end{proposition}
\begin{proof}
The onsite sum in~\eqref{eq:a25} vanishes, the edge sum is $4D$, and
$d_{\rm pre}=(-d_0+g(q-1),-2g(q-1),d_0+g(q-1))$.
Opposite onsite terms cancel the summed mixed term; the full square is
$2d_0^2+12g^2D$. Substitution gives~\eqref{eq:a27}. The synchronizer has not been suppressed: its modulus preservation explains why $\lambda$ does not appear.

The alternating cosine series, whose term magnitudes decrease on the stated interval, gives
\begin{equation}\label{eq:a28}% accepted equation 28
\begin{gathered}
\frac{\phi^2}{2}-\frac{\phi^4}{24}\le1-\cos\phi\le\frac{\phi^2}{2}.
\end{gathered}
\end{equation}
At $\phi=1/10000$ this yields $\Delta I\ge387/20000000000>0$.
At $\phi=1/1000$ it gives $\Delta I\le-107999987/2400000000000000<0$.
Since $D$ is continuous and strictly increasing for these positive phases, the unique threshold is
\begin{equation}\label{eq:a29}% accepted equation 29
\begin{gathered}
\phi_*=\arccos(1-1/6500000),\qquad 1/10000<\phi_*<1/1000.
\end{gathered}
\end{equation}
The baseline is positive below, zero at, and negative above this threshold.
\end{proof}
Dropping the finite square would lose the positive constant and change the coefficient of $D$. Against the old positive launch ceiling $C_\phi=\alpha D$,
\begin{equation}\label{eq:a30}% accepted equation 30
\begin{gathered}
\frac{b_0}{C_\phi}=\frac1{100000000D}-\frac{13}{200}\longrightarrow+\infty
\quad(\phi\downarrow0).
\end{gathered}
\end{equation}
The numerator tends to the finite value $\alpha/100000000$, while the denominator tends to zero. This is not infinite accessible storage or an efficiency claim.

\subsection{Departure from the earlier exact preparation}\label{sec:familyexit}
Direct expansion gives
\begin{equation}\label{eq:a31}% accepted equation 31
\begin{gathered}
|F_1\Omega_\phi|^2=(1-d_0)^2-2g(1-d_0-g)D<(1-d_0)^2<1,
\end{gathered}
\end{equation}
because $D>0$ and $1-d_0-g>0$. Also
\begin{equation}\label{eq:a32}% accepted equation 32
\begin{gathered}
|F_3\Omega_\phi|^2-|F_1\Omega_\phi|^2\\ =4d_0(1-gD)>0.
\end{gathered}
\end{equation}
Here $D\le1/20000$, so $1-gD>0$. Every allowed positive phase therefore leaves the exact unit-modulus family. Unequal first/third moduli cannot be repaired by one common phase or common nonzero amplitude factor. In particular $F\Omega_\phi\ne\Omega_\phi$.

This is a failure of the narrow launch hypothesis of the earlier receiving/recovery theorem, not a proof of departure from the larger native coherent domain or impossibility of hybrid receiving. Even an interval $h$ exceeding the old return time only leaves time before the next jump: in the fixed native-first protocol the first exchange starts at $F\Omega_\phi$, not $\Omega_\phi$. Subsequent intervals also generally retain external preload.

\section{Actual post-native source and initial receiving}\label{sec:activity}
\subsection{Phase order through simultaneous synchronization}\label{sec:phaseorder}
Write $v_i=\rho_i e^{i\theta_i}$. The prestage has the positive-cone representation
\begin{equation}\label{eq:a33}% accepted equation 33
\begin{gathered}
v_1=(1-d_0-g)+gq,\\ v_2=2g+(1-2g)q,\\ v_3=(1+d_0-g)+gq.
\end{gathered}
\end{equation}
All coefficients are positive. For $0<\phi\le1/100$, each $v_i$ is nonzero and strictly between the rays of $1$ and $q$, so $0<\theta_i<\phi$ with no wrap. Direct multiplication gives
\begin{equation}\label{eq:a34}% accepted equation 34
\begin{gathered}
\operatorname{Im}(\overline v_1v_2)=A_*\sin\phi,\\ A_*=1-3g-d_0(1-2g)=\frac{284971}{300000}>0.
\end{gathered}
\end{equation}
Likewise $\operatorname{Im}(\overline v_3v_1)=2d_0g\sin\phi>0$. All phase differences have magnitude below $\pi/2$, so sine has their sign and $\theta_3<\theta_1<\theta_2$.

For the actual updated phases $\theta_i^+=\theta_i+\delta_i$, use the following weights and bound for $|u|\le\phi$:
\begin{equation}\label{eq:a35}% accepted equation 35
\begin{gathered}
w(u)=\begin{cases}\lambda\sin(3u)/u,&u\ne0,\\3\lambda,&u=0.\end{cases}
\quad 0\le w(u)\le3\lambda .
\end{gathered}
\end{equation}
The inequality uses $|3u|\le3/100<\pi$ and $|\sin(3u)|\le3|u|$. Then
\begin{equation}\label{eq:a36}% accepted equation 36
\begin{gathered}
\theta_i^+=\left(1-\sum_{j\ne i}w(\theta_j-\theta_i)\right)\theta_i+
\sum_{j\ne i}w(\theta_j-\theta_i)\theta_j.
\end{gathered}
\end{equation}
The off-diagonal weight sum is at most $6\lambda=1/50<1$. Every updated phase remains between the original minimum and maximum, inside $(0,\phi)$. These are simultaneous weights, not a sequential averaging scheme.

For $\theta_j>\theta_i$, write $\Delta=\theta_j-\theta_i>0$ and let $\ell$ be the third node. Subtracting phase increments gives
\begin{equation}\label{eq:a37}% accepted equation 37
\begin{gathered}
\theta_j^+-\theta_i^+
=\Delta-2\lambda\sin(3\Delta)
\\ {}+\lambda[\sin(3(\theta_\ell-\theta_j))-\sin(3(\theta_\ell-\theta_i))]
\ge(1-9\lambda)\Delta.
\end{gathered}
\end{equation}
Indeed $\sin(3\Delta)\le3\Delta$ and the derivative of $\sin(3u)$ has magnitude at most $3$, so the bracket is at least $-3\Delta$. With $s_*=1-9\lambda=97/100>0$, all strict orderings survive.

In particular $\Delta^+=\theta_2^+-\theta_1^+\ge s_*\Delta>0$ and $\Delta^+\le\phi<\pi/2$. Monotonicity and concavity of sine imply $\sin\Delta^+\ge\sin(s_*\Delta)\ge s_*\sin\Delta$, whence
\begin{equation}\label{eq:a38}% accepted equation 38
\begin{gathered}
\chi_3^+=\rho_1\rho_2\sin\Delta^+
\ge s_*\rho_1\rho_2\sin\Delta=s_*A_*\sin\phi>0.
\end{gathered}
\end{equation}

\subsection{The registered source and the actual local result}\label{sec:source}
The raw port convention remains
\begin{equation}\label{eq:a39}% accepted equation 39
\begin{gathered}
f_{{\rm raw},\pm}=\frac G{\sqrt7}
\{\sqrt2(\chi_1+\chi_2)\ \mathord{\pm}\ \sqrt3\chi_3\},\\ f_{{\rm raw},+}-f_{{\rm raw},-}=\frac{2\sqrt3G}{\sqrt7}\chi_3.
\end{gathered}
\end{equation}
At the original launch $\chi=(-\sin\phi,0,\sin\phi)$ this gives $G\sin\phi(-\sqrt2\pm\sqrt3)/\sqrt7$, with positive raw $+$ and negative raw $-$ entries. For the post-native result only the strictly positive difference in~\eqref{eq:a39} is needed; no individual post-native port sign is assumed. Squaring and adding gives
\begin{equation}\label{eq:a40}% accepted equation 40
\begin{gathered}
|K\chi|^2=\frac{G^2}{7}\{4(\chi_1+\chi_2)^2+6\chi_3^2\}.
\end{gathered}
\end{equation}
Together with~\eqref{eq:a38}, this proves
\begin{equation}\label{eq:a41}% accepted equation 41
\begin{gathered}
|f(F\Omega_\phi)|^2\ge
\frac{6G^2}{7}s_*^2A_*^2\sin^2\phi>0.
\end{gathered}
\end{equation}
The exact coefficient multiplying $G^2\sin^2\phi$ is
$764090502142969/1050000000000000$.
Uniformity concerns that coefficient, not an absolute positive floor as $\phi\downarrow0$.

\begin{theorem}[Initial supply and receiving after the actual native jump]\label{thm:initial}
Fix one finite positive exchange tuple before phase. For each fixed $0<\phi\le1/100$, launch the actual flow at $(F\Omega_\phi,0,0,0)$ and put $f_{\rm post}=f(F\Omega_\phi)$. There is a $\delta(\phi,\mathrm{tuple})>0$ such that native exchange debit and receiver gain are strict for $0<t<\min(h,\delta)$. Their leading terms, and the receiver state, are given by~\eqref{eq:a43}; both leading interface powers are positive as in~\eqref{eq:a44}. No phase-uniform duration or later receiving/recovery episode is asserted.
\end{theorem}
\begin{proof}
The general empty-carrier proposition~\cite[Proposition 4.1]{rx} applies to arbitrary native preparation with nonzero drive. Equation~\eqref{eq:a41} supplies that hypothesis at the actual synchronized state; no special RR rate relation is imposed. Empty carrier data belong to $H^1$, so the strong exchange theorem gives $p'(0)=\eta\B f_{\rm post}$ and $r_c'(0)=0$ in $L^2$. Boundedness of $\B^*$ gives
\begin{equation}\label{eq:a42}% accepted equation 42
\begin{gathered}
p(t)=\eta tBf_{\rm post}+o_{L^2}(t),\\ r_c(t)=o_{L^2}(t),\\ y(t)=\eta mtf_{\rm post}+o(t).
\end{gathered}
\end{equation}
Local boundedness of $J$ and $y(t)=O(t)$ imply $x(t)-x(0)=O(t^2)$, and the quadratic $f$ gives $f(x(t))=f_{\rm post}+O(t^2)$. Integration of the exact native and receiver equations yields
\begin{equation}\label{eq:a43}% accepted equation 43
\begin{aligned}
H_N(F\Omega_\phi)-H_N(t)&=\frac{\eta^2m}{2}|f_{\rm post}|^2t^2+o(t^2),\\
z(t)&=\frac{\eta\kappa m}{2\beta}f_{\rm post}t^2+o(t^2),\\
H_R(t)&=\frac{\eta^2\kappa^2m^2}{8\beta}|f_{\rm post}|^2t^4+o(t^4).
\end{aligned}
\end{equation}
The receiver vector expansion is squared in its original store $\beta|z|^2/2$. More explicitly, the signed interface powers satisfy
\begin{equation}\label{eq:a44}% accepted equation 44
\begin{gathered}
P_{NC}(t)=\eta^2m|f_{\rm post}|^2t+o(t),\\ P_{CR}(t)=\frac{\eta^2\kappa^2m^2}{2\beta}|f_{\rm post}|^2t^3+o(t^3).
\end{gathered}
\end{equation}
All leading coefficients are positive. At each fixed preparation, the little-$o$ remainders can be bounded by half the corresponding leading coefficients on a sufficiently short interval. Taking the minimum of those finitely many neighbourhood sizes proves strict positivity. Moreover $Q_{N\to C}>0$ there, because the sum of positive per-port supply powers is at least the positive signed total. The next native jump truncates the interval at $h$.
\end{proof}
Receiver feedback is retained throughout; this is not a clamped-source comparison. The argument uses first $L^2$ derivatives and integration, not a second derivative of a tent. The coefficient estimate~\eqref{eq:a41} alone supplies no phase-uniform relative remainder. No conclusion at the next sample $h$ follows unless it lies in the short-time regime. In particular the later quarter-capacity receiving and greater-than-$24/25$ native return results remain unproved for this hybrid.

\section{A local comparison of algebraic capacities}\label{sec:capacity}
\subsection{A different active branch near zero phase}\label{sec:branch}
For the original launch $I_0=3$, $Z_0=2+e^{2i\phi}$, $h_e=D$ and $j_0=1+2\cos\phi$. Both compared quantities are positive and
\begin{equation}\label{eq:a45}% accepted equation 45
\begin{gathered}
j_0^2-|Z_0|^2=4\cos\phi(1-\cos\phi)>0.
\end{gathered}
\end{equation}
so the original ceiling uses the projection branch, $C_\phi=\alpha D$~\cite[Section 3, strict yaw witness]{pc}.

Keep the actual detuning $d_0\ne0$ fixed. At zero phase,
\begin{equation}\label{eq:a46}% accepted equation 46
\begin{gathered}
\mathbf r_0=F\Omega_0=(1-d_0,1,1+d_0),\\ R^2=|\mathbf r_0|^2=3+2d_0^2,\\ C_Z(\mathbf r_0)=0,\quad h_e(\mathbf r_0)=d_0^2/2>0.
\end{gathered}
\end{equation}
The continuous difference $\alpha h_e-C_Z$ is strictly positive there, hence stays positive in some neighbourhood of zero. Thus the post-native joint capacity uses $C_Z$ in that neighbourhood. No explicit width or full-interval branch claim follows from this argument.

\subsection{Derivative, parity and the finite remainder}\label{sec:capacityderivative}
Let $u=(g,1-2g,g)$. The prestage derivative at zero is $iu$. Since every component of $\mathbf r_0$ is positive, local phase extraction is real analytic and $\theta_i'(0)=u_i/(\mathbf r_0)_i$. Differentiating the simultaneous synchronizer gives
\begin{equation}\label{eq:a47}% accepted equation 47
\begin{gathered}
\delta_i'(0)=3\lambda\left(\sum_j u_j/(\mathbf r_0)_j-3u_i/(\mathbf r_0)_i\right),\\ F\Omega_\phi=\mathbf r_0+i\xi\phi+O(\phi^2),\\ \xi=s_*u+3\lambda \mathbf r_0\sum_j u_j/(\mathbf r_0)_j .
\end{gathered}
\end{equation}
Equivalently $\xi=\operatorname{diag}(\mathbf r_0)(\mathrm{Id}+3\lambda L_3)\operatorname{diag}(1/\mathbf r_0)u$. The radial term contributes no cross product:
\begin{equation}\label{eq:a48}% accepted equation 48
\begin{gathered}
\mathbf r_0\times \xi=s_*(\mathbf r_0\times u)
=\frac1{30000000}(-27647813,97,27642187).
\end{gathered}
\end{equation}
Locally $F\Omega_{-\phi}=\overline{F\Omega_\phi}$: the prestage has real coefficients and the phases and sine increments are odd under conjugation near the positive limiting state. Analyticity and parity therefore give
\begin{equation}\label{eq:a49}% accepted equation 49
\begin{gathered}
\operatorname{Re}F\Omega_\phi=\mathbf r_0+\zeta\phi^2+O(\phi^4),\\ \operatorname{Im}F\Omega_\phi=\xi\phi+O(\phi^3),\\ \chi^+=(\mathbf r_0\times \xi)\phi+O(\phi^3).
\end{gathered}
\end{equation}
The unknown real coefficient $\zeta$ must not be discarded. Its contributions are
\begin{equation}\label{eq:a50}% accepted equation 50
\begin{aligned}
I&=R^2+(2\mathbf r_0\cdot \zeta+|\xi|^2)\phi^2+O(\phi^4),\\
|Z_2|&=R^2+
\left(2\mathbf r_0\cdot \zeta-|\xi|^2+\frac{2(\mathbf r_0\cdot \xi)^2}{R^2}\right)\phi^2+O(\phi^4).
\end{aligned}
\end{equation}
The modulus expansion is valid because $Z_2(0)=R^2>0$. Subtraction cancels $\zeta$:
\begin{equation}\label{eq:a51}% accepted equation 51
\begin{gathered}
C_Z(F\Omega_\phi)=\alpha\left(|\xi|^2-\frac{(\mathbf r_0\cdot \xi)^2}{R^2}\right)\phi^2
+O(\phi^4).
\end{gathered}
\end{equation}
Equivalently, $C_Z=2\alpha|a\times b|^2/(I+|Z_2|)$ and the even nonzero denominator give the same remainder order. The fixed $\alpha$ is included in the constant implicit in $O(\phi^4)$.

The radial term in $\xi$ cancels from its squared transverse norm. Exact arithmetic yields
\begin{equation}\label{eq:a52}% accepted equation 52
\begin{gathered}
|u|^2=\frac{187}{200},\quad \mathbf r_0\cdot u=1,\\ q_{\rm cap}=s_*^2\left(\frac{187}{200}-\frac1{3+2d_0^2}\right)
=\frac{169832451759483}{300000002000000}>0.
\end{gathered}
\end{equation}
Combining the strict branch selection with this expansion and $1-\cos\phi=\phi^2/2+O(\phi^4)$ proves
\begin{equation}\label{eq:a53}% accepted equation 53
\begin{gathered}
C(F\Omega_\phi)=\alpha q_{\rm cap}\phi^2+O(\phi^4),\\ \lim_{\phi\downarrow0}\frac{C(F\Omega_\phi)}{C_\phi}
=2q_{\rm cap}=\frac{169832451759483}{150000001000000}
\simeq1.132216337515111.
\end{gathered}
\end{equation}
This is a local comparison of algebraic ceilings, not an achieved transfer fraction or efficiency improvement. Sending detuning to zero first removes the strict branch gap and changes the argument. No simultaneous zero-detuning limit, explicit phase neighbourhood or detuning optimization is supplied.

\section{Small-cadence limits and their scope}\label{sec:cadence}
\begin{proposition}[One fixed state]\label{prop:onestep}
For every fixed $X\in\mathcal H$,
\begin{equation}\label{eq:a54}% accepted equation 54
\begin{gathered}
T_hX=\Phi_h(\widehat FX)\longrightarrow\widehat FX
\quad\mathrm{in}\ \|\cdot\|_E,\quad h\downarrow0.
\end{gathered}
\end{equation}
If $\widehat FX\ne X$, then
\begin{equation}\label{eq:a55}% accepted equation 55
\begin{gathered}
\|T_hX-X\|_E\longrightarrow\|\widehat FX-X\|_E>0,\\ \frac{\|T_hX-X\|_E}{h}\longrightarrow+\infty.
\end{gathered}
\end{equation}
\end{proposition}
\begin{proof}
The image $\widehat FX$ is a defined finite-energy point. Strong continuity of exchange at that fixed point proves~\eqref{eq:a54}, without requiring continuity of $F$ at neighbouring inputs. Norm continuity gives the positive limit of the numerator in~\eqref{eq:a55}. No differentiation of the unbounded translation generator is used.
\end{proof}
Equation~\eqref{eq:a31} makes every permitted positive-phase launch nonfixed. In fact $|F_1-1|\ge1-|F_1|>d_0$. If samples at time $h$ approximated uniformly near zero a continuous curve $Y$ with $Y(0)=X$, then $Y(h)\to X$ while $T_hX\to\widehat FX\ne X$. The first-sample error could not tend to zero. Thus fixed finite native updates do not give ordinary near-identity continuous refinement with that same initial value. Useful finite cadences, special fixed states and separately formulated singular or initial-layer limits are not excluded.

\begin{proposition}[Fixed finite number of updates]\label{prop:finiteN}
Fix a finite $N$. Assume $F$ is continuous at the native parts of $X,\widehat FX,\ldots,\widehat F^{N-1}X$. Then~\eqref{eq:a57} holds. Nonzero prestage components at these states are sufficient.
\end{proposition}
\begin{proof}
Away from zero prestages, powers of $v_i/|v_i|$ express the phase factors smoothly; an argument branch cut alone is no obstruction. At zero-prestage strata the convention defines $F$ but supplies no general neighbouring continuity.

The inherited bounded-energy dependence estimate~\cite[Theorem 3.1]{rx} gives, for nearby $Y_h\to Y$ and $h\downarrow0$,
\begin{equation}\label{eq:a56}% accepted equation 56
\begin{gathered}
\|\Phi_hY_h-Y\|_E
\le e^{Lh}\|Y_h-Y\|_E+\|\Phi_hY-Y\|_E\longrightarrow0
\quad\text{if }Y_h\to Y,\ h\downarrow0,
\end{gathered}
\end{equation}
The constant $L$ in this estimate is a local Lipschitz constant, not the invariant floor $\mathcal L$. Continuity of $\widehat F$ at the encountered point and this joint continuity of exchange advance a finite induction. Therefore
\begin{equation}\label{eq:a57}% accepted equation 57
\begin{gathered}
T_h^NX\longrightarrow\widehat F^NX\qquad\text{for each such fixed finite }N.
\end{gathered}
\end{equation}
\end{proof}
The inherited native-only forward-invariant coherent-domain result supplies nonzero prestages along native-only coherent orbits~\cite[Section 6, native domain]{rx}. This suffices for any fixed finite $N$ there, but does not assert invariance of an active hybrid trajectory. No constants uniform in $N$ are proved, so substituting $N$ proportional to $1/h$ is not justified.

\section{A dark boundary with positive native baseline}\label{sec:dark}
The aligned boundary $\phi=0$ is outside the strictly positive receiving family, but is an admissible separate native preparation. Exactly,
\begin{equation}\label{eq:a58}% accepted equation 58
\begin{gathered}
F(1,1,1)=(9999/10000,1,10001/10000)=\mathbf r_0,\\ b=\frac\alpha{100000000}>0.
\end{gathered}
\end{equation}
Before and after the jump the state is real, so $\chi=f=0$, $Z_2=I$, $\mathcal L=H_N$ and $C=0$. In detail,
\begin{equation}\label{eq:a59}% accepted equation 59
\begin{aligned}
I(\mathbf r_0)=Z_2(\mathbf r_0)&=3+2d_0^2,&
h_e(\mathbf r_0)&=d_0^2/2,& j(\mathbf r_0)&=3+d_0^2,\\
\Delta I=\Delta Z_2&=2d_0^2,&
\Delta j&=d_0^2,& \Delta \mathcal L&=b,\quad\Delta C=0.
\end{aligned}
\end{equation}
Both exchange invariants change, while the whole baseline increase raises the invariant floor and leaves capacity unchanged.

More generally, a real native state has a real polynomial prestage. Its arguments are integer multiples of $\pi$, with the specified zero convention, so every sine in the mathematical synchronizer vanishes. The image remains real even at zero-prestage components. This is an exact statement, not bitwise floating-point equality for $\sin(3\pi)$.

With empty external coordinates a real native state is stationary under exchange: $f=y=z=0$ makes the constant full state a solution, and uniqueness applies. Induction gives real native states and zero external stores at every finite hybrid iterate, for any fixed positive $h$. Later baseline increments need not all be positive; the first increment above is the witness. With arbitrary external preload, stationarity is no longer guaranteed, because $J^Ty$ can be nonzero~\cite[Section 3, zero ceiling and drive]{pc}.

The boundary witness is consistent with nonzero drive for every strictly positive phase and with capacity tending to zero. It shows that native baseline gain alone guarantees neither activity nor capacity replenishment.

\section{Interpretation, evidence and open questions}\label{sec:discussion}
\subsection{What remains free}\label{sec:open}
The first actual post-native exchange segment has strict initial native debit and receiver gain at each fixed allowed positive phase and finite positive coefficient tuple, only on a sufficiently short interval before the next jump. Theorem~\ref{thm:initial} resolves initial activity; it does not establish the later capacity-sized receiving-and-recovery episode under native stepping. Later rounds also carry actual external preload. Baseline-corrected identities remain exact, but the endpoint estimates needed for a quantified return have not been proved for those windows.

The comparison rates $\nu,\gamma,\omega_R$ and the earlier exchange-only times permit dimensionless comparisons such as $\nu h,\gamma h,\omega_Rh$ and $ch/L_c$~\cite[Section 3]{rr}. They do not select $h$ or identify native iteration count with physical time. A source-off carrier period describes actual propagation only on its source-off hypothesis; an autonomous active feedback system is not source-off merely because it has no external forcing~\cite[Section 6]{rx}.

No phase-uniform positive short-time interval, explicit local capacity-branch width, growing-$N$ limit, uniform infinite-step bound, preferred cadence or physical clock is proved. These are limits of the present result, not mandatory repairs to its accepted claims. Failure to inherit the old episode theorem is not impossibility of hybrid receiving or return. Exchange-first ordering, delayed first jumps, state resets, normalization, a hidden compensating reservoir or coefficients scaled with $h$ would specify different problems.

\subsection{Provenance and evidence}\label{sec:evidence}
This paper is prepared from the accepted F01--F12 investigation and its actual GPT review~\cite{investigation,acceptance}. The eight exact-algebra groups support the written arguments; the two original checker replays and the later GPT replay do not multiply independent findings. Regularity, endpoint conventions, strict inequalities, branch selection, parity and remainders are analytical arguments. No scientific replay was needed for typesetting, and no new trajectory, phase sweep or cadence optimization was run.

The accompanying crosswalk maps every accepted equation, claim/limitation row, upstream source locator and algebra group to final manuscript labels and pages. The accepted sources and historical ledgers remain unchanged. The two stale internal-note path references are resolved only in new mappings. Retained source records are not described as fresh remote retrievals.

GPT formulated and internally reviewed the investigation; Codex independently investigated it and prepared this manuscript under the research author's direction. Scientific acceptance does not certify this new PDF: the manuscript and proposed publication payload await GPT/author review. No external peer review or machine-formal proof is claimed.

\subsection{Geometric and research boundary}\label{sec:continuity}
The energetic account does not reconstruct a spatial circuit, assign a material receiver location or supply irreversible retention or physical calibration. The admitted carrier circle, intended tight toroid, close original sphere and extended-line sphere remain distinct; finite host-notch N regions are distinct from the three upper and three lower central G boundaries.

The continuing programme retains the central starting/meeting region, opposed cones, outward spreading, curved around-and-inward return, stationary supports and possible transfer elsewhere as geometric intentions. Intended central-route correspondence and then six-opening transport remain subsequent directions, with gravity--magnetism linked throughout. Outer Tri-Octagon geometry, distinct golden-ratio relationships and justified SRG integration stand later; no new ordering or scientific premise is supplied here. Three-Way remains parked; completed simulation reconstructions and the bridge remain resources. Pi and H\'ok\'i p\'ok\'i remain brainstorming.

A physical-time identification would require an additional stated timing rule. The present results supply no such calibration.

\clearpage
\begin{thebibliography}{9}
\bibitem{rx}
H.~Fr\'imann Halld\'orsson,
\emph{A Reciprocal Chiral Exchange Extension of the Tri-Octagon Model},
v0.1 (2026).
\href{https://github.com/pzychozen/trioctagon-physics/tree/273b1a6faaab19ea58748d00576439166eeef9ab/papers/RECIPROCAL_CHIRAL_EXCHANGE/publication/v0.1}{Immutable edition and evidence}.
Theorem 3.1: global exchange flow; Proposition 4.1: empty-carrier activity; Section 7: transfer accounts; Section 6: exact native map and hybrid.
\bibitem{pc}
H.~Fr\'imann Halld\'orsson,
\emph{Recoverable Transfer and Projection Capacity in the Reciprocal Chiral Exchange Candidate},
v0.1 (2026).
\href{https://github.com/pzychozen/trioctagon-physics/tree/54c0f6c419fa681e2c1e74a5ad2bb46df787bb93/papers/RECOVERABLE_TRANSFER_AND_PROJECTION_CAPACITY/publication/v0.1}{Immutable edition and evidence}.
Theorem 3.1: general projection invariant and joint ceiling; Section 3: algebraic sharpness and its dynamical limits.
\bibitem{rr}
H.~Fr\'imann Halld\'orsson,
\emph{Capacity-Sized Receiving and Native Recovery in the Reciprocal Chiral Exchange Candidate},
v0.1 (2026).
\href{https://github.com/pzychozen/trioctagon-physics/tree/34bf6c6e4e18adcc177ac700359fda16cf4d772d/papers/NATIVE_CAPACITY_RECEIVING_RECOVERY/publication/v0.1}{Immutable scientific edition and evidence}.
Theorem 3.1 is the continuous-exchange episode certificate for its exact coherent empty launch.
\bibitem{nativecode}
H.~Fr\'imann Halld\'orsson,
\emph{Native dynamics definitions}, \texttt{kernel\_physics/dynamics.py},
repository snapshot 6b1ef6494a4d4c34398ba9d9e5a7afce1fb03f37 (2026).
\href{https://github.com/pzychozen/trioctagon-physics/blob/6b1ef6494a4d4c34398ba9d9e5a7afce1fb03f37/kernel_physics/dynamics.py}{Immutable source}.
Definitions \texttt{DynamicsConfig}, \texttt{arg0}, \texttt{phase\_sync}, \texttt{\_advance} and \texttt{step3}; read as mathematical definitions, not executed for this manuscript.
\bibitem{investigation}
H.~Fr\'imann Halld\'orsson, with GPT formulation and Codex independent investigation,
\emph{Native F-stepping and exchange cadence: independent proof review},
9 November 2026. Accepted obligations F01--F12, equations 1--59, claim/source ledger and eight exact-algebra support groups, preserved in the accompanying evidence.
\bibitem{acceptance}
GPT, \emph{Native F stepping and exchange cadence: GPT acceptance},
9 November 2026. Project-internal scientific acceptance with the stated hypotheses; no mathematical correction required. Complete acceptance, audits and historical replay are preserved unchanged. This decision does not approve the new typeset manuscript.
\end{thebibliography}

\end{document}
